% !PRINTED 513
\weberpage{513}

We shall now understand by $x_0, y_0, \dots$ any integral rational numbers whatever. Then
\[ \delta_0 = \alpha x_0 + \beta y_0 + \cdots \]
is also an integral functional, and we shall prove \gesperrt{that $\delta_0$ is divisible by $\delta$.}

For this purpose we need only consider the integral functional
\[ \delta t - \delta_0 = \alpha (x t - x_0) + \beta (y t - y_0) \dots \]
in which $t$ denotes a new variable. Then [by (2)]:

\weq{3}{N (\delta t - \delta_0) = D E (x t - x_0, y t - y_0, \dots),}

or, if we arrange according to descending powers of $t$:

\weq{4}{N (\delta t - \delta_0) = D (t^n E + t^{n-1} E_1 + t^{n-2} E_2 + \cdots),}

in which $E_1, E_2, \dots$ are, by \S.\,137, 7., 8., integral rational functionals. If we now set
\[ C_1 = \frac{E_1}{E}, \quad C_2 = \frac{E_2}{E}, \cdots, \]
then, since $E$ is a unit, $C_1, C_2, \dots$ are also integral rational functionals; and if we further set
\[ \frac{\delta_0}{\delta} = \eta \]
it follows from (4) that
\[ N (\delta) N (t - \eta) = D E (t^n + C_1 t^{n-1} + C_2 t^{n-2} + \cdots), \]
or, on account of (2),

\weq{5}{N (t - \eta) = t^n + C_1 t^{n-1} + C_2 t^{n-2} + \cdots}

Now since this function vanishes when we set $t = \eta$, it follows that $\eta$ is an integral functional (\S.\,137, 2.), and it is thus proved that $\delta_0$ is divisible by $\delta$.

Since $x_0, y_0, \dots$ may denote arbitrary integral rational numbers, we conclude from this that the functionals $\alpha, \beta, \dots$ are themselves divisible by $\delta$, and hence that $\delta$ is a common divisor of the functionals $\alpha, \beta, \dots$.

On the other hand, however, $\delta$ is also (by \S.\,138, 4.) divisible by every common divisor of $\alpha, \beta, \dots$, and $\delta$ therefore has the \gesperrt{characteristic properties of the greatest common divisor of $\alpha, \beta, \dots$} The same properties, however, belong by \S.\,138, 7. to every functional associated with $\delta$; and likewise two functionals $\delta, \delta'$ that have this double prop-
%CONT>>
